% GENERATED FILE. Edit canonical lessons, then run npm run generate:books.
% canonical-source-hash: fnv1a64:65ad9e9d17fb3f3c
% canonical-lessons: TA-C359-tiri, TA-C359-mikaippatuttu, TA-C359-muranpatu, TA-C359-aiyam, TA-C359-nirakari

\chapter{Twisting the Truth}
\label{ch:ta-twisting-the-truth}
\hlchaptermodality{359}

\begin{chapteropening}
\textbf{By the end of this chapter:} \emph{I can say that someone twists, exaggerates or rejects what was said.}

Complete the last lesson of chapter 359: I can say that someone twists, exaggerates or rejects what was said.
\end{chapteropening}

\section[\texorpdfstring{\ta{திரி} (tiri)}{திரி (tiri)}]{\ta{திரி} (tiri) — to twist; to distort}

\label{lesson:TA-C359-tiri}



\begin{quote}
Before the new one: say the Tamil for a policy, then the Tamil for to establish.
\end{quote}

\subsection*{You'll want to know: \ta{திரி}}
\textbf{\ta{திரி}} — \emph{tiri} — \textquotedblleft{}to twist; to distort\textquotedblright{}.

\begin{quote}
\ta{அவர்} \ta{என்} \ta{வார்த்தைகளைத்} \ta{திரித்துச்} \ta{சொன்னார்}.
\end{quote}

\begin{quote}
\emph{avar eṉ vārttaikaḷait tirittuc coṉṉār.}
\end{quote}

\begin{quote}
He twisted my words.
\end{quote}

Past (I …): \ta{திரித்தேன்}, \emph{tirittēṉ}.

Transitive with -\ta{ஐ}, past \ta{திரித்தேன்}; the intransitive \ta{திரி}, to wander, has \ta{திரிந்தேன்}.

\begin{cousinweb}
\textbf{\ta{திரி}} (\emph{tiri}) is native Dravidian (DEDR 3246): with the past in -\ta{ந்த}- it means \textquotedblleft{}to turn, wander, change\textquotedblright{}; with the past in -\ta{த்த}- it means \textquotedblleft{}to twist\textquotedblright{}, as a rope is twisted.
\end{cousinweb}

\begin{grammarlens}[title={What you've built}]
One more word for giving reasons and opinions.
\end{grammarlens}

\subsection*{Your turn}
\begin{itemize}
\raggedright
  \item \emph{Say it:} \emph{tiri}
  \item \emph{Say it:} \emph{tiri}, once more
  \item \emph{Say it:} \emph{avar eṉ vārttaikaḷait tirittuc coṉṉār.} — He twisted my words.
  \item \emph{Recall:} say the Tamil for a policy, then the Tamil for to establish, then say \emph{tiri} again
\end{itemize}

\begin{tcolorbox}[breakable,colback=teal!4,colframe=teal!35!black,title={Before you move on}]
What does \ta{திரி} mean? (\textquotedblleft{}To twist; to distort\textquotedblright{}.) Say it once more.
\end{tcolorbox}
\section[\texorpdfstring{\ta{மிகைப்படுத்து} (mikaippaṭuttu)}{மிகைப்படுத்து (mikaippaṭuttu)}]{\ta{மிகைப்படுத்து} (mikaippaṭuttu) — to exaggerate}

\label{lesson:TA-C359-mikaippatuttu}



\begin{quote}
Before the new one: say the Tamil for to establish, then the Tamil for to twist.
\end{quote}

\subsection*{You'll want to know: \ta{மிகைப்படுத்து}}
\textbf{\ta{மிகைப்படுத்து}} — \emph{mikaippaṭuttu} — \textquotedblleft{}to exaggerate\textquotedblright{}.

\begin{quote}
\ta{அவன்} \ta{எப்போதும்} \ta{எல்லாவற்றையும்} \ta{மிகைப்படுத்துகிறான்}.
\end{quote}

\begin{quote}
\emph{avaṉ eppōtum ellāvaṟṟaiyum mikaippaṭuttukiṟāṉ.}
\end{quote}

\begin{quote}
He always exaggerates everything.
\end{quote}

Past (I …): \ta{மிகைப்படுத்தினேன்}, \emph{mikaippaṭuttiṉēṉ}.

\ta{மிகை}, excess: \ta{மிகையாகச்} \ta{சொல்}, to overstate.

\begin{cousinweb}
\textbf{\ta{மிகைப்படுத்து}} (\emph{mikaippaṭuttu}) is native: \emph{mikai}, \textquotedblleft{}abundance, excess\textquotedblright{} (DEDR 4838), with \emph{paṭuttu}, the causative of \emph{paṭu}: \textquotedblleft{}to make more than it is\textquotedblright{}.
\end{cousinweb}

\begin{grammarlens}[title={What you've built}]
One more word for giving reasons and opinions.
\end{grammarlens}

\subsection*{Your turn}
\begin{itemize}
\raggedright
  \item \emph{Say it:} \emph{mikaippaṭuttu}
  \item \emph{Say it:} \emph{mikaippaṭuttu}, once more
  \item \emph{Say it:} \emph{avaṉ eppōtum ellāvaṟṟaiyum mikaippaṭuttukiṟāṉ.} — He always exaggerates everything.
  \item \emph{Recall:} say the Tamil for to establish, then the Tamil for to twist, then say \emph{mikaippaṭuttu} again
\end{itemize}

\begin{tcolorbox}[breakable,colback=teal!4,colframe=teal!35!black,title={Before you move on}]
What does \ta{மிகைப்படுத்து} mean? (\textquotedblleft{}To exaggerate\textquotedblright{}.) Say it once more.
\end{tcolorbox}
\section[\texorpdfstring{\ta{முரண்படு} (muraṇpaṭu)}{முரண்படு (muraṇpaṭu)}]{\ta{முரண்படு} (muraṇpaṭu) — to contradict; to clash}

\label{lesson:TA-C359-muranpatu}



\begin{quote}
Before the new one: say the Tamil for to twist, then the Tamil for to exaggerate.
\end{quote}

\subsection*{You'll want to know: \ta{முரண்படு}}
\textbf{\ta{முரண்படு}} — \emph{muraṇpaṭu} — \textquotedblleft{}to contradict; to clash\textquotedblright{}.

\begin{quote}
\ta{அவர்} \ta{சொன்னது} \ta{ஆதாரத்துடன்} \ta{முரண்படுகிறது}.
\end{quote}

\begin{quote}
\emph{avar coṉṉatu ātārattuṭaṉ muraṇpaṭukiṟatu.}
\end{quote}

\begin{quote}
What he said contradicts the evidence.
\end{quote}

Past (I …): \ta{முரண்பட்டேன்}, \emph{muraṇpaṭṭēṉ}.

With -\ta{உடன்}: X-\ta{உடன்} \ta{முரண்படு}. The noun is \ta{முரண்பாடு}, a contradiction.

\begin{cousinweb}
\textbf{\ta{முரண்படு}} (\emph{muraṇpaṭu}) is native: DEDR 4971 lists Tamil \emph{muraṇ}, \textquotedblleft{}to be at variance, be opposed; variance\textquotedblright{}, here with \emph{paṭu}, \textquotedblleft{}to undergo\textquotedblright{}.
\end{cousinweb}

\begin{grammarlens}[title={What you've built}]
One more word for giving reasons and opinions.
\end{grammarlens}

\subsection*{Your turn}
\begin{itemize}
\raggedright
  \item \emph{Say it:} \emph{muraṇpaṭu}
  \item \emph{Say it:} \emph{muraṇpaṭu}, once more
  \item \emph{Say it:} \emph{avar coṉṉatu ātārattuṭaṉ muraṇpaṭukiṟatu.} — What he said contradicts the evidence.
  \item \emph{Recall:} say the Tamil for to twist, then the Tamil for to exaggerate, then say \emph{muraṇpaṭu} again
\end{itemize}

\begin{tcolorbox}[breakable,colback=teal!4,colframe=teal!35!black,title={Before you move on}]
What does \ta{முரண்படு} mean? (\textquotedblleft{}To contradict; to clash\textquotedblright{}.) Say it once more.
\end{tcolorbox}
\section[\texorpdfstring{\ta{ஐயம்} (aiyam)}{ஐயம் (aiyam)}]{\ta{ஐயம்} (aiyam) — uncertainty; a doubt}

\label{lesson:TA-C359-aiyam}



\begin{quote}
Before the new one: say the Tamil for to exaggerate, then the Tamil for to contradict.
\end{quote}

\subsection*{You'll want to know: \ta{ஐயம்}}
\textbf{\ta{ஐயம்}} — \emph{aiyam} — \textquotedblleft{}uncertainty; a doubt\textquotedblright{}.

\begin{quote}
\ta{இதில்} \ta{எனக்கு} \ta{எந்த} \ta{ஐயமும்} \ta{இல்லை}.
\end{quote}

\begin{quote}
\emph{itil eṉakku enta aiyamum illai.}
\end{quote}

\begin{quote}
I have no doubt about this.
\end{quote}

A formal word; the everyday one is \ta{சந்தேகம்}. \ta{ஐயமின்றி}, without doubt.

\begin{cousinweb}
\textbf{\ta{ஐயம்}} (\emph{aiyam}) is native Dravidian: DEDR 190 lists Tamil \emph{ayir}, \textquotedblleft{}to suspect\textquotedblright{}, and \emph{aiyam}, \textquotedblleft{}doubt, uncertainty\textquotedblright{}, with Malayalam \emph{ayyam}, \textquotedblleft{}doubt\textquotedblright{}.
\end{cousinweb}

\begin{grammarlens}[title={What you've built}]
One more word for giving reasons and opinions.
\end{grammarlens}

\subsection*{Your turn}
\begin{itemize}
\raggedright
  \item \emph{Say it:} \emph{aiyam}
  \item \emph{Say it:} \emph{aiyam}, once more
  \item \emph{Say it:} \emph{itil eṉakku enta aiyamum illai.} — I have no doubt about this.
  \item \emph{Recall:} say the Tamil for to exaggerate, then the Tamil for to contradict, then say \emph{aiyam} again
\end{itemize}

\begin{tcolorbox}[breakable,colback=teal!4,colframe=teal!35!black,title={Before you move on}]
What does \ta{ஐயம்} mean? (\textquotedblleft{}Uncertainty; a doubt\textquotedblright{}.) Say it once more.
\end{tcolorbox}
\section[\texorpdfstring{\ta{நிராகரி} (nirākari)}{நிராகரி (nirākari)}]{\ta{நிராகரி} (nirākari) — to reject; to turn down}

\label{lesson:TA-C359-nirakari}



\begin{quote}
Before the new one: say the Tamil for to contradict, then the Tamil for uncertainty.
\end{quote}

\subsection*{You'll want to know: \ta{நிராகரி}}
\textbf{\ta{நிராகரி}} — \emph{nirākari} — \textquotedblleft{}to reject; to turn down\textquotedblright{}.

\begin{quote}
\ta{நிறுவனம்} \ta{அவருடைய} \ta{விண்ணப்பத்தை} \ta{நிராகரித்தது}.
\end{quote}

\begin{quote}
\emph{niṟuvaṉam avaruṭaiya viṇṇappattai nirākarittatu.}
\end{quote}

\begin{quote}
The company rejected his application.
\end{quote}

Past (I …): \ta{நிராகரித்தேன்}, \emph{nirākarittēṉ}.

What is rejected takes -\ta{ஐ}. A formal word; \ta{மறு} is the everyday refuse.

\begin{cousinweb}
\textbf{\ta{நிராகரி}} (\emph{nirākari}) is a verb built on Sanskrit \emph{nirākaraṇa}, \textquotedblleft{}driving away, repudiating, denying\textquotedblright{} (Monier-Williams), with the loan-verb ending -\ta{இ}.
\end{cousinweb}

\begin{grammarlens}[title={What you've built}]
One more word for giving reasons and opinions.
\end{grammarlens}

\subsection*{Your turn}
\begin{itemize}
\raggedright
  \item \emph{Say it:} \emph{nirākari}
  \item \emph{Say it:} \emph{nirākari}, once more
  \item \emph{Say it:} \emph{niṟuvaṉam avaruṭaiya viṇṇappattai nirākarittatu.} — The company rejected his application.
  \item \emph{Recall:} say the Tamil for to contradict, then the Tamil for uncertainty, then say \emph{nirākari} again
\end{itemize}

\begin{tcolorbox}[breakable,colback=teal!4,colframe=teal!35!black,title={Before you move on}]
What does \ta{நிராகரி} mean? (\textquotedblleft{}To reject; to turn down\textquotedblright{}.) Say it once more.
\end{tcolorbox}
